%----------------------------------------------------------------------------------------
%	ABBREVIATIONS and GLOSSARY
%----------------------------------------------------------------------------------------

% --- ChatGPT and aliases ---
\newglossaryentry{chatgpt}{
    name={ChatGPT},
    description={An AI language model developed by OpenAI that can understand and generate human-like text based on prompts}
}
\newglossaryentry{chat-gpt}{
    name={Chat-GPT},
    description={See \gls{chatgpt}},
    see=[main entry:]{chatgpt}
}
\newglossaryentry{Chat-GPT}{
    name={Chat GPT},
    description={See \gls{chatgpt}},
    see=[main entry:]{chatgpt}
}

% --- GPT-4o ---
\newacronym{gpt4o}{GPT-4o}{Generative Pre-trained Transformer 4 Omni}
\newglossaryentry{gpt-4o}{
    name={GPT-4o},
    description={A multimodal \gls{large-language-model} developed by OpenAI that processes text, vision, and audio inputs natively and in real time},
    see=[acronym:]{gpt4o}
}

% --- DeepSeek-R1 ---
\newacronym{deepseekr1}{DeepSeek-R1}{DeepSeek Research Model 1}
\newglossaryentry{deepseek-r1}{
    name={DeepSeek-R1},
    description={An open-source \gls{large-language-model} developed by DeepSeek, designed for general-purpose reasoning and code understanding},
    see=[acronym:]{deepseekr1}
}

% --- Python Programming Language ---
\newglossaryentry{python}{
    name={Python},
    description={A high-level, interpreted programming language known for its readability, simplicity, and versatility in fields such as web development, data science, automation, and artificial intelligence}
}
\newglossaryentry{python-language}{
    name={Python language},
    description={See \gls{python}},
    see=[main entry:]{python}
}

% --- API and aliases ---
\newacronym{api}{API}{Application Programming Interface}
\newglossaryentry{application-programming-interface}{
    name={Application Programming Interface},
    description={A set of protocols and tools that allows different software applications to communicate and interact with each other},
    see=[acronym:]{api}
}
\newglossaryentry{api-entry}{
    name={API},
    description={See \gls{application-programming-interface}},
    see=[main entry:]{application-programming-interface}
}

% --- RAG and aliases ---
\newacronym{rag}{RAG}{Retrieval-Augmented Generation}
\newglossaryentry{retrieval-augmented-generation}{
    name={Retrieval-Augmented Generation},
    description={A technique that combines information retrieval with language generation to improve the accuracy and relevance of AI-generated responses},
    see=[acronym:]{rag}
}
\newglossaryentry{rag-approach}{
    name={RAG approach},
    description={See \gls{retrieval-augmented-generation}},
    see=[main entry:]{retrieval-augmented-generation}
}
\newglossaryentry{rag-method}{
    name={RAG method},
    description={See \gls{retrieval-augmented-generation}},
    see=[main entry:]{retrieval-augmented-generation}
}

% --- LLM and aliases ---
\newacronym{llm}{LLM}{Large Language Model}
\newglossaryentry{large-language-model}{
    name={Large Language Model},
    description={A type of AI trained on vast text data to understand, generate, and manipulate human language},
    see=[acronym:]{llm}
}
\newglossaryentry{language-model}{
    name={language model},
    description={See \gls{large-language-model}},
    see=[main entry:]{large-language-model}
}
\newglossaryentry{LLMs}{
    name={LLMs},
    description={See \gls{large-language-model}},
    see=[main entry:]{large-language-model}
}

% --- AI and aliases ---
\newacronym{ai}{AI}{Artificial Intelligence}
\newglossaryentry{artificial-intelligence}{
    name={Artificial Intelligence},
    description={The simulation of human intelligence in machines designed to think, learn, and solve problems},
    see=[acronym:]{ai}
}
\newglossaryentry{machine-intelligence}{
    name={machine intelligence},
    description={See \gls{artificial-intelligence}},
    see=[main entry:]{artificial-intelligence}
}
\newglossaryentry{ai-systems}{
    name={AI systems},
    description={See \gls{artificial-intelligence}},
    see=[main entry:]{artificial-intelligence}
}

% --- Static Malware Analysis and aliases ---
\newglossaryentry{static-malware-analysis}{
    name={Static Malware Analysis},
    description={The process of examining malware without executing it, typically by analyzing its code, structure, and metadata to identify its functionality and potential threats}
}
\newglossaryentry{static-analysis}{
    name={static analysis},
    description={See \gls{static-malware-analysis}},
    see=[main entry:]{static-malware-analysis}
}

% --- Machine Learning and aliases ---
\newacronym{ml}{ML}{Machine Learning}
\newglossaryentry{machine-learning}{
    name={Machine Learning},
    description={A subset of artificial intelligence that enables systems to learn and improve from experience without being explicitly programmed},
    see=[acronym:]{ml}
}
\newglossaryentry{ml-system}{
    name={ML system},
    description={See \gls{machine-learning}},
    see=[main entry:]{machine-learning}
}
\newglossaryentry{ml-model}{
    name={ML model},
    description={See \gls{machine-learning}},
    see=[main entry:]{machine-learning}
}

% --- Malware and aliases ---
\newglossaryentry{malware}{
    name={Malware},
    description={Malicious software designed to disrupt, damage, or gain unauthorized access to computer systems}
}
\newglossaryentry{malicious-software}{
    name={malicious software},
    description={See \gls{malware}},
    see=[main entry:]{malware}
}
\newglossaryentry{malware-sample}{
    name={malware sample},
    description={See \gls{malware}},
    see=[main entry:]{malware}
}

% --- Decompilation and aliases ---
\newglossaryentry{decompilation}{
    name={Decompilation},
    description={The process of translating low-level machine code or bytecode back into a higher-level, human-readable source code}
}
\newglossaryentry{decompile}{
    name={decompile},
    description={See \gls{decompilation}},
    see=[main entry:]{decompilation}
}
\newglossaryentry{decompiles}{
    name={decompiles},
    description={See \gls{decompilation}},
    see=[main entry:]{decompilation}
}
\newglossaryentry{decompiler}{
    name={Decompiler},
    description={See \gls{decompilation}},
    see=[main entry:]{decompilation}
}

% --- PostgreSQL and aliases ---
\newglossaryentry{postgresql}{
    name={PostgreSQL},
    description={An advanced, open-source relational database management system known for its extensibility and SQL compliance}
}
\newglossaryentry{postgres}{
    name={Postgres},
    description={See \gls{postgresql}},
    see=[main entry:]{postgresql}
}

% --- pgvector and aliases ---
\newglossaryentry{pgvector}{
    name={pgvector},
    description={An extension for \gls{postgresql} that provides support for vector similarity search, enabling efficient handling of embeddings and machine learning workloads}
}
\newglossaryentry{pg-vector}{
    name={pg-vector},
    description={See \gls{pgvector}},
    see=[main entry:]{pgvector}
}

% --- Ghidra and aliases ---
\newglossaryentry{ghidra}{
    name={Ghidra},
    description={An open-source software reverse engineering (SRE) suite developed by the NSA, used for analyzing compiled code, decompilation, and malware analysis}
}
\newglossaryentry{ghydra}{
    name={Ghydra},
    description={See \gls{ghidra}},
    see=[main entry:]{ghidra}
}

% --- Haystack (Python library) and aliases ---
\newglossaryentry{haystack}{
    name={Haystack (Python library)},
    description={An open-source Python framework for building large language model (LLM) applications, including question answering, retrieval-augmented generation (RAG), and document search}
}
\newglossaryentry{haystack-llm}{
    name={Haystack LLM},
    description={See \gls{haystack}},
    see=[main entry:]{haystack}
}
\newglossaryentry{haystack-python}{
    name={Haystack Python library},
    description={See \gls{haystack}},
    see=[main entry:]{haystack}
}

% --- peframe and aliases ---
\newglossaryentry{peframe}{
    name={PEFrame},
    description={An open-source tool used for static analysis of Portable Executable (PE) files to detect malware and extract useful information}
}
\newglossaryentry{pe-frame}{
    name={PE frame},
    description={See \gls{peframe}},
    see=[main entry:]{peframe}
}

% --- PE file and aliases ---
\newglossaryentry{pe-file}{
    name={PE file},
    description={A Portable Executable (PE) file, the standard file format for executables, object code, and DLLs in 32-bit and 64-bit versions of Windows operating systems}
}
\newglossaryentry{portable-executable}{
    name={Portable Executable},
    description={See \gls{pe-file}},
    see=[main entry:]{pe-file}
}
\newglossaryentry{pe-files}{
    name={PE files},
    description={See \gls{pe-file}},
    see=[main entry:]{pe-file}
}
\newglossaryentry{pe}{
    name={PE format},
    description={The file format specification for Portable Executable (PE) files used in Windows operating systems; defines the structure and layout of executable files, object code, and DLLs}
}

% --- VirusTotal ---
\newglossaryentry{virustotal}{
    name={VirusTotal},
    description={A free online service that analyses files and URLs for viruses, worms, trojans, and other kinds of malicious content using multiple antivirus engines}
}

% --- MalwareBazaar ---
\newglossaryentry{malwarebazaar}{
    name={MalwareBazaar},
    description={A platform for sharing and downloading malware samples, aimed at aiding cybersecurity researchers and analysts in studying malware}
}

% --- Haystacks---
\newglossaryentry{haystacks}{
    name={Haystacks},
    description={An open-source AI framework designed for building end-to-end pipelines for tasks such as question answering, document retrieval, and natural language processing}
}
